\section{Gaussian Space --- Product Hessian}
%%% generated by the blueprint pipeline from TCSlib.BooleanAnalysis.GaussianSpace.ProductHessian
\subsection*{Gaussian Space --- Product Hessian: Introduction}
The mean score and Hessian of a product Gaussian density have explicit linear and quadratic coordinate formulas. Local integrable envelopes justify a second differentiation of bounded Gaussian averages and yield global Lipschitz derivative bounds.

\subsection*{Gaussian Space --- Product Hessian: Results}

\begin{definition}[Product Gaussian linear score]
\lean{BooleanAnalysis.GaussianSpace.gaussianProductDensityScore}
\leanok
Let $I$ be finite and $v\ge0$. Write $p_v(c,z)$ for the $N(c,v)$ density when $v>0$ and set $p_0(c,z)=0$; define $G_v(c,z)=\prod_{i\in I}p_v(c_i,z_i)$ and interpret division by zero as zero. For $c,z,h\in\mathbb R^I$, put $S_v(c,z)[h]=\sum_i((z_i-c_i)/v)h_i$. The mean score is the linear functional $S_v(c,z)$.
\end{definition}

\begin{lemma}[Derivative of the product Gaussian score]
\lean{BooleanAnalysis.GaussianSpace.hasFDerivAt_gaussianProductDensityScore}
\leanok
\uses{BooleanAnalysis.GaussianSpace.gaussianProductDensityScore}
\difficulty{4}
Let $I$ be a finite index set, let $v\ge0$, and let $c,z\in\mathbb R^I$. The function
$m\mapsto S_v(m,z)$, where $S_v(m,z)[h]=\sum_{i\in I}((z_i-m_i)/v)h_i$, is
differentiable at $c$. Its derivative sends a direction $h$ to the linear functional
$k\mapsto-\sum_{i\in I}h_i k_i/v$, with division by zero interpreted as zero.
\end{lemma}

\begin{definition}[Product Gaussian mean Hessian]
\lean{BooleanAnalysis.GaussianSpace.gaussianProductDensityHessian}
\leanok
\uses{BooleanAnalysis.GaussianSpace.gaussianProductDensityDerivative, BooleanAnalysis.GaussianSpace.gaussianProductDensityScore}
Let $I$ be finite and $v\ge0$. Put $G_v(c,z)=\prod_i p_v(c_i,z_i)$, where $p_v$ is the Gaussian density of variance $v$ for $v>0$ and $p_0=0$, and put $S_v(c,z)[h]=\sum_i((z_i-c_i)/v)h_i$, with division by zero equal to zero. Define $A_v(c,z)[h]=G_v(c,z)S_v(c,z)[h]$ and $T_v(c,z)[h,k]=G_v(c,z)(S_v(c,z)[h]S_v(c,z)[k]-\sum_i h_i k_i/v)$. The product-density Hessian is the bilinear form $T_v(c,z)$.
\end{definition}

\begin{lemma}[Derivative of the product density derivative]
\lean{BooleanAnalysis.GaussianSpace.hasFDerivAt_gaussianProductDensityDerivative}
\leanok
\uses{BooleanAnalysis.GaussianSpace.gaussianProductDensityDerivative, BooleanAnalysis.GaussianSpace.gaussianProductDensityDerivative_eq_density_smul_score, BooleanAnalysis.GaussianSpace.gaussianProductDensityHessian, BooleanAnalysis.GaussianSpace.gaussianProductDensityScore, BooleanAnalysis.GaussianSpace.hasFDerivAt_gaussianProductDensity, BooleanAnalysis.GaussianSpace.hasFDerivAt_gaussianProductDensityScore}
\difficulty{3}
Let $I$ be a finite index set, let $v\ge0$, and let $c,z\in\mathbb R^I$. The function of
$m\in\mathbb R^I$ given by the derivative of the product Gaussian density at mean $m$ is
differentiable at $c$. Its derivative, applied to directions $h,k\in\mathbb R^I$, is
$G_v(c,z)\bigl(S_v(c,z)[h]S_v(c,z)[k]-\sum_{i\in I}h_i k_i/v\bigr)$, where $G_v$ is the
product Gaussian density, $S_v(c,z)[h]=\sum_i((z_i-c_i)/v)h_i$, and division by zero is
interpreted as zero.
\end{lemma}

\begin{lemma}[Coordinate formula for the Gaussian Hessian]
\lean{BooleanAnalysis.GaussianSpace.gaussianProductDensityHessian_eq_sum}
\leanok
\uses{BooleanAnalysis.GaussianSpace.gaussianProductDensityDerivative_eq_density_smul_score, BooleanAnalysis.GaussianSpace.gaussianProductDensityHessian, BooleanAnalysis.GaussianSpace.gaussianProductDensityScore}
\difficulty{5}
Let $I$ be finite, $v\ge0$, and $c,z\in\mathbb R^I$. The product Gaussian density
Hessian, as a continuous bilinear form, is the sum over $i,j\in I$ of the rank-one forms
$(h,k)\mapsto h_jk_i$, each multiplied by
$\left((z_i-c_i)(z_j-c_j)/v^2-\delta_{ij}/v\right)\prod_{k\in I}p_v(c_k,z_k)$, where
$p_v(a,b)$ denotes the one-dimensional Gaussian density with mean $a$ and variance $v$,
and $\delta_{ij}$ is $1$ if $i=j$ and $0$ otherwise.
\end{lemma}

\begin{lemma}[Bound for a squared Gaussian score]
\lean{BooleanAnalysis.GaussianSpace.gaussianProductDensity_squaredScore_le}
\leanok
\difficulty{5}
Let $I$ be finite, let $v\ge0$, and let $c,z,a\in\mathbb R^I$ satisfy $a_i\ge0$,
$p_v(c_i,z_i)\le a_i$, and $\left|\left((z_i-c_i)^2/v^2-1/v\right)p_v(c_i,z_i)\right|\le
a_i$ for every $i\in I$, where $p_v(u,w)$ denotes the Gaussian density with mean $u$ and
variance $v$. Then for each $i\in I$, $\frac{(z_i-c_i)^2}{v^2}\prod_{j\in
I}p_v(c_j,z_j)\le(1+1/v)\prod_{j\in I}a_j$.
\end{lemma}

\begin{lemma}[Bound for a product Gaussian Hessian coefficient]
\lean{BooleanAnalysis.GaussianSpace.abs_gaussianProductDensity_hessianCoefficient_le}
\leanok
\uses{BooleanAnalysis.GaussianSpace.abs_productHessian_mul_le, BooleanAnalysis.GaussianSpace.gaussianProductDensity_squaredScore_le}
\difficulty{5}
Let $I$ be a finite index set, let $v\ge0$ be a real number with $v\ne0$, and let
$c,z,a:I\to\mathbb R$. Suppose $a(i)\ge0$ and, for every $i\in I$, the Gaussian density
with mean $c(i)$ and variance $v$, evaluated at $z(i)$, is at most $a(i)$, and
\[
\left|\left(\frac{(z(i)-c(i))^2}{v^2}-\frac1v\right)p_v(c(i),z(i))\right|\le a(i),
\]
where $p_v$ denotes that density. Then for every $i,j\in I$,
\[
\left|\left(\frac{(z(i)-c(i))(z(j)-c(j))}{v^2}-\begin{cases}1/v,&i=j,\\0,&i\ne
j\end{cases}\right)\prod_{k\in I}p_v(c(k),z(k))\right|\le\left(1+\frac2v\right)\prod_{k
i I}a(k).
\]
\end{lemma}

\begin{lemma}[Joint continuity of the product density Hessian]
\lean{BooleanAnalysis.GaussianSpace.continuous_gaussianProductDensityHessian_joint}
\leanok
\uses{BooleanAnalysis.GaussianSpace.gaussianProductDensityHessian, BooleanAnalysis.GaussianSpace.gaussianProductDensityHessian_eq_sum}
\difficulty{4}
Let $I$ be a finite set and let $v\ge 0$. The product Gaussian density Hessian, viewed
as a continuous bilinear form on $\mathbb R^I$, is jointly continuous as a function of
its mean vector $c\in\mathbb R^I$ and input vector $z\in\mathbb R^I$.
\end{lemma}

\begin{lemma}[Norm bound for a coordinate tensor sum]
\lean{BooleanAnalysis.GaussianSpace.norm_sum_sum_coordinateTensor_le}
\leanok
\uses{BooleanAnalysis.GaussianSpace.norm_coordinateProjection_le_one}
\difficulty{4}
Let $I$ be a finite set and let $a_{ij}\in\mathbb R$ for $i,j\in I$. Define the bilinear
form $T\colon\mathbb R^I\times\mathbb R^I\to\mathbb R$ by $T(u,v)=\sum_{i,j\in
I}a_{ij}u_jv_i$, and equip $\mathbb R^I$ with its usual supremum norm. Its operator norm
satisfies $\|T\|\leq\sum_{i,j\in I}|a_{ij}|$.
\end{lemma}

\begin{lemma}[Integrable envelope for the Gaussian Hessian]
\lean{BooleanAnalysis.GaussianSpace.gaussianProductDensityHessian_dominated}
\leanok
\uses{BooleanAnalysis.GaussianSpace.abs_gaussianProductDensity_hessianCoefficient_le, BooleanAnalysis.GaussianSpace.gaussianPDFReal_two_derivatives_dominated, BooleanAnalysis.GaussianSpace.gaussianProductDensityHessian, BooleanAnalysis.GaussianSpace.gaussianProductDensityHessian_eq_sum, BooleanAnalysis.GaussianSpace.norm_sum_sum_coordinateTensor_le}
\difficulty{5}
Let $I$ be finite, let $v\ge0$ with $v\ne0$, and let $M\ge0$. There is a nonnegative
function $b:\mathbb R^I\to\mathbb R$ that is integrable with respect to Lebesgue measure
on $\mathbb R^I$ and bounds the operator norm of the product Gaussian density Hessian at
every $z\in\mathbb R^I$, uniformly over all $m\in\mathbb R^I$ with $|m_i|\le M$ for
every $i\in I$.
\end{lemma}

\begin{lemma}[Integrability of the weighted Gaussian density Hessian]
\lean{BooleanAnalysis.GaussianSpace.integrable_gaussianProductDensityHessian_smul}
\leanok
\uses{BooleanAnalysis.GaussianSpace.continuous_gaussianProductDensityHessian_joint, BooleanAnalysis.GaussianSpace.gaussianProductDensityHessian, BooleanAnalysis.GaussianSpace.gaussianProductDensityHessian_dominated}
\difficulty{4}
Let $I$ be a finite set, let $f:\mathbb R^I\to\mathbb R$ be measurable with $|f(z)|\le
1$ for every $z$, let $v>0$, and let $c\in\mathbb R^I$. Define $G_v(c,z)=\prod_i
p_v(c_i,z_i)$ and $S_v(c,z)[h]=\sum_i((z_i-c_i)/v)h_i$. The function sending $z$ to
$f(z)$ times the bilinear form $(h,k)\mapsto G_v(c,z)\bigl(S_v(c,z)[h]S_v(c,z)[k]-\sum_i
h_i k_i/v\bigr)$ is integrable with respect to Lebesgue measure on $\mathbb R^I$, as a
bilinear-form-valued function.
\end{lemma}

\begin{lemma}[Derivative of the integrated density derivative]
\lean{BooleanAnalysis.GaussianSpace.hasFDerivAt_integral_gaussianProductDensityDerivative}
\leanok
\uses{BooleanAnalysis.GaussianSpace.abs_coordinate_le_of_mem_ball, BooleanAnalysis.GaussianSpace.continuous_gaussianProductDensityDerivative, BooleanAnalysis.GaussianSpace.continuous_gaussianProductDensityHessian_joint, BooleanAnalysis.GaussianSpace.gaussianProductDensityDerivative, BooleanAnalysis.GaussianSpace.gaussianProductDensityHessian, BooleanAnalysis.GaussianSpace.gaussianProductDensityHessian_dominated, BooleanAnalysis.GaussianSpace.hasFDerivAt_gaussianProductDensityDerivative, BooleanAnalysis.GaussianSpace.integrable_gaussianProductDensityDerivative_smul}
\difficulty{5}
Let $I$ be a finite set, let $v>0$, and let $f\colon\mathbb R^I\to\mathbb R$ be
measurable with $|f(z)|\le1$ for every $z$. The function $m\mapsto\int_{\mathbb
R^I}f(z)A_v(m,z)\,dz$, where $A_v$ is the mean derivative of the product Gaussian
density of variance $v$, is Fréchet differentiable at every $c\in\mathbb R^I$, with
derivative $h\mapsto\int_{\mathbb R^I}f(z)T_v(c,z)[h,\cdot]$; here $T_v$ is the Hessian
of the product density, and the displayed derivative is the linear map obtained by
integrating the bilinear form in its second argument.
\end{lemma}

\begin{lemma}[Integral of the Gaussian product density Hessian]
\lean{BooleanAnalysis.GaussianSpace.integral_gaussianProductDensityHessian_smul_eq_sum}
\leanok
\uses{BooleanAnalysis.GaussianSpace.gaussianProductDensityHessian, BooleanAnalysis.GaussianSpace.gaussianProductDensityHessian_eq_sum, BooleanAnalysis.GaussianSpace.integrable_productHessian_mul, BooleanAnalysis.GaussianSpace.integral_pi_gaussianReal_eq_integral_smul}
\difficulty{5}
Let $I$ be a finite index set, let $v>0$, let $c\in\mathbb R^I$, and let $f:\mathbb
R^I\to\mathbb R$ be measurable with $|f(z)|\le1$ for every $z$. Let $Z$ have independent
coordinates with $Z_k\sim N(c_k,v)$. The integral, with respect to Lebesgue measure on
$\mathbb R^I$, of $f(z)$ times the product Gaussian mean Hessian equals
\[\sum_{i,j\in I}\mathbb E\!\left[\left(\frac{(Z_i-c_i)(Z_j-c_j)}{v^2}-\frac{\mathbf 1_{i=j}}{v}\right)f(Z)\right] (e_j^*\otimes e_i^*),\]
where $e_k^*$ is the coordinate projection $z\mapsto z_k$, and $e_j^*\otimes e_i^*$
denotes the bilinear form $(h,k)\mapsto h_jk_i$.
\end{lemma}

\begin{lemma}[Bound for an integrated Gaussian density Hessian]
\lean{BooleanAnalysis.GaussianSpace.norm_integral_gaussianProductDensityHessian_smul_le}
\leanok
\uses{BooleanAnalysis.GaussianSpace.abs_integral_productHessian_mul_le, BooleanAnalysis.GaussianSpace.gaussianProductDensityHessian, BooleanAnalysis.GaussianSpace.integral_gaussianProductDensityHessian_smul_eq_sum, BooleanAnalysis.GaussianSpace.norm_sum_sum_coordinateTensor_le}
\difficulty{4}
Let $I$ be a finite index set, and let $f:\mathbb R^I\to\mathbb R$ be measurable with
$|f(z)|\leq1$ for every $z$. Let $v\geq0$ with $v\ne0$, and let $c\in\mathbb R^I$. The
norm of the integral of $f(z)$ times the product Gaussian density Hessian with mean $c$
and variance $v$, integrated against Lebesgue measure on $\mathbb R^I$, is at most
\[
|I|^2\frac{2}{v}.
\]
\end{lemma}

\begin{lemma}[Lipschitz continuity of Gaussian product density averages]
\lean{BooleanAnalysis.GaussianSpace.lipschitzWith_integral_gaussianProductDensityDerivative}
\leanok
\uses{BooleanAnalysis.GaussianSpace.gaussianProductDensityDerivative, BooleanAnalysis.GaussianSpace.hasFDerivAt_integral_gaussianProductDensityDerivative, BooleanAnalysis.GaussianSpace.norm_integral_gaussianProductDensityHessian_smul_le}
\difficulty{3}
Let $I$ be a finite index set, let $f:\mathbb R^I\to\mathbb R$ be measurable with
$|f(z)|\le1$ for every $z$, and let $v>0$. For mean vector $c$, integrate $f(z)$ times
the mean derivative of the product Gaussian density with variance $v$ in each coordinate
over Lebesgue measure on $\mathbb R^I$. As a function of $c$, this
linear-functional-valued integral is Lipschitz with constant $2|I|^2/v$, using the
maximum norm on mean vectors and the operator norm on linear functionals.
\end{lemma}
